\subsection*{2.4 极端暗光环境下融合偏振特征点的视觉自主导航}

室内、地下、夜间室外等极端暗光环境对视觉导航系统的可靠性与鲁棒性提出了严峻挑战。在照度低于数十lux的暗光条件下，常规可见光相机获取的图像信噪比极低，传统视觉特征点提取与跟踪方法面临特征数量骤减、误匹配率激增的问题，导致视觉惯性里程计（VIO）的几何约束严重不足，甚至完全失效。直接法视觉里程计依赖光度不变假设，在暗光高噪声条件下该假设难以成立，同样无法有效工作。

针对上述问题，本节以微阵列偏振相机与IMU为传感器配置，提出一种极端暗光环境下融合偏振特征点的视觉自主导航方法。其核心思路为：利用偏振信息（偏振度与偏振角）在暗光条件下仍能捕捉环境结构信息的物理优势，从微阵列偏振相机原始数据中提取多维度偏振通道图像；经专门设计的预处理流程抑制暗光噪声后，在各偏振通道上提取视觉特征点；通过跨通道全局特征池将多通道特征按空间位置关联融合，弥补可见光通道的特征丢失；最终在滑动窗口非线性优化框架下实现偏振-可见光融合的视觉惯性导航，使系统在极端暗光条件下仍能维持稳定、精确的自主导航能力。

如1.4.2节所述，现有偏振SLAM工作或服务于稠密建图、或提供后端航向约束、或面向户外低纹理场景，尚无工作针对极端暗光条件下偏振通道的前端特征增强能力进行系统性研究。本节正是针对这一空白展开，提出以下三项核心创新点：

\textbf{(a) 基于偏振通道提取与导向滤波预处理的暗光特征增强方法。}针对暗光条件下可见光图像信噪比极低、传统特征提取方法失效的问题，从微阵列偏振相机的原始偏振数据中解算偏振度（DoP）与偏振角（AoP）的$\sin$/$\cos$分量，构造三个偏振通道图像$P_0, P_1, P_2$；系统对比中值滤波、导向滤波、双边滤波与非局部均值去噪四种预处理算法在不同光照条件下的特征跟踪性能，确立以$S_0$强度图像为引导的导向滤波作为暗光偏振通道的最优预处理策略。该方法在抑制暗光椒盐噪声的同时保留边缘结构信息，使极暗条件下（0.9--4.8 lux）的特征匹配内点比例由无滤波的0.244--0.360提升至0.549--0.650，为暗光视觉特征跟踪提供高质量输入。

\textbf{(b) 基于跨通道全局特征池的偏振-可见光特征融合策略。}针对暗光下$S_0$可见光通道特征大量丢失而偏振通道$P_0, P_1, P_2$仍保有结构信息的特点，设计跨通道全局特征池（Global Feature Pool）机制：以像素空间位置为关联准则，将四个通道独立检测的局部特征按传播-关联-输出三阶段流程合并为统一的全局特征，以固定通道优先级（$S_0 \rightarrow P_0 \rightarrow P_1 \rightarrow P_2$）向后端输出观测。该策略使特征的生命周期不再依赖单一通道——只要至少一个关联通道仍在跟踪，全局特征即可持续存在，从而在暗光条件下显著增加可跟踪特征数量与有效观测数：极暗条件下（0.9--4.8 lux）融合全部偏振通道的匹配特征数约为仅$S_0$通道的6倍，有效观测数提升约2.5倍。

\textbf{(c) 暗光偏振融合视觉惯性导航系统（PolarFP-VINS）与实验验证。}将上述偏振特征提取与跨通道融合方法集成至VINS-Fusion后端优化框架，构建完整的PolarFP-VINS系统；在涵盖94 lux至0.9 lux的8组真实室内暗光数据集上进行系统级验证，与VINS-Fusion基线方法及DSO、D3VO等SOTA直接法进行全面对比。实验结果表明，在最大光强20 lux及以下的暗光条件下，PolarFP-VINS的终端位置误差（TPE）显著优于基线方法；在极暗条件（0.9--4.8 lux）下，基线方法TPE达2.70 m而PolarFP-VINS仅为1.19 m；DSO与D3VO在所有暗光数据上均无法完成轨迹解算，验证了偏振信息在暗光视觉导航中的独特价值与不可替代性。

\subsubsection*{2.4.1 暗光条件下的偏振特征提取方法}

\textbf{(1) 偏振通道提取}

为提取偏振特征，首先要从微阵列偏振相机的原始数据中提取偏振通道信息。

具体来说，微阵列式偏振相机的偏振感知原理，是在图像传感器上覆盖了微透镜与微偏振片阵列，如图\ref{fig:micro_polar_array}所示。在微阵列式偏振相机中，每四个相邻像素组成一组超像素。一组超像素基本对应同一空间方向射来的光，可以认为对应普通相机的一个像素。该组超像素中的每一子像素，分别感知该方向的光，经过相对偏振角零位$0^\circ$/$45^\circ$/$90^\circ$/$135^\circ$偏振片后，产生的偏振光强$I_0/I_{45}/I_{90}/I_{135}$。偏振角的零位方向是相机感光元件的横轴右方向，即相机坐标系的$x$轴方向。

\begin{figure}[H]
\centering
\includegraphics[width=0.9\linewidth]{pic/1_polar_dark/1_micro_polar_array.png}
\caption{微阵列偏振相机像素排布(左)与结构示意图(右)}
\label{fig:micro_polar_array}
\end{figure}

利用同一超像素的四个偏振光强，可以由公式~\ref{eq:stokes} 恢复该像素的斯托克斯矢量。其中，斯托克斯矢量第一维$S_0$即为总光强，也就是普通光学相机捕捉到的光强。

\begin{equation}
\label{eq:stokes}
\begin{aligned}
S_0 &= \frac{1}{2}\bigl(I_0 + I_{45} + I_{90} + I_{135}\bigr), \\
S_1 &= I_0 - I_{90}, \\
S_2 &= I_{45} - I_{135}, \\
S_3 &= 0.
\end{aligned}
\end{equation}

公式~\ref{eq:stokes} 中，斯托克斯矢量的最后一维是圆偏振光分量。由于自然界中圆偏振光占比较少，且现有微偏振阵列式偏振相机不支持原偏振光感知，在本研究中，我们假设圆偏振光分量均可忽略，即斯托克斯矢量$S_3$均为0。此时，线偏振度(DolP, Degree of linear Polarization)与总偏振度(DoP, Degree of Polarization)相等.

基于上述假设，根据公式~\ref{eq:aop-and-dop}，我们可以从斯托克斯矢量计算偏振度与偏振角(AoP, Angle of Polarization)。

\begin{equation}
\label{eq:aop-and-dop}
\begin{aligned}
DoP &= \frac{\sqrt{S_1^2 + S_2^2}}{S_0}, \\
AoP &= \frac{1}{2} \arctan\frac{S_2}{S_1}.
\end{aligned}
\end{equation}


值得注意的是，在极端暗光环境下，$S_0$可能非常接近零，直接按公式~\ref{eq:aop-and-dop}计算$DoP$会出现除零不稳定；此外，图像噪声可能导致$S_1,S_2$的模长异常大于$S_0$，出现$DoP>1$的非物理情况。因此，在实际实现中，我们对分母加入小常数$\varepsilon=10^{-6}$防止除零；同时对$S_0<\varepsilon$或$DoP>0.999$的像素，将其$DoP$、$\sin(AoP)$、$\cos(AoP)$均置零，以消除无效区域的异常偏振响应。

对图像传感器的每组超像素均进行上述计算，我们就得到偏振度和偏振角组成的图像。上述计算得到的偏振度取值范围为$[0,1]$，偏振角取值为$[0, \pi]$。为便于进行数字图像处理，我们对偏振角计算其$\sin$和$\cos$值，和偏振度一起组成3个偏振通道$p_0, p_1, p_2$，如公式~\ref{eq:polar-channel}所示

\begin{equation}
\label{eq:polar-channel}
\begin{aligned}
p_0 & = DoP, \\
p_1 & = \sin(AoP), \\
p_2 & = \cos(AoP) \\
\end{aligned}
\end{equation}

对于这三个偏振通道，我们通过线性变换将像素值取值范围映射到$[0, 255]$，并取整量化为相邻整数值，得到8位数字图像$P_0, P_1, P_2$。其中，$P_0$直接由$[0,1]$线性放大255倍；$P_1$和$P_2$的原始范围为$[-1,1]$，先放大127.5倍再进行127.5的亮度平移，从而完整映射到$[0,255]$。在后文中，我们称的“偏振通道”均为$P_0, P_1, P_2$而非$p_0, p_1, p_2$。

\begin{equation}
\label{eq:polar-channel-2}
\begin{aligned}
P_0 & = \left[p_0 \times 255 \right], \\
P_1 & = \left[p_1 \times \frac{255}{2} + \frac{255}{2}\right] = \left[p_1 \times 127.5 + 127.5\right], \\
P_2 & = \left[p_2 \times \frac{255}{2} + \frac{255}{2}\right] = \left[p_2 \times 127.5 + 127.5\right] \\
\end{aligned}
\end{equation}

在对应正常光照和不同暗弱光照条件的三种不同光强下，对同一场景采集3组偏振图像，作为可视化展示数据。具体来说，第1组数据采集时室内光强为94-205 lux(勒克斯)，对应正常室内光照条件；第2组数据采集时室内光强为2.6-18.6 lux，对应较暗室内光照条件；第3组数据采集时室内光强为0.9-4.8 lux，对应极暗室内光照条件。

使用以上方法，对各光照条件下采集的原始偏振相机图像进行处理，得到的斯托克斯矢量、各偏振通道如图\ref{fig:polar_channels}所示。在该图中，斯托克斯矢量$S_1, S_2$的可视化，是将原始值取绝对值后，进行30倍增益后显示的结果。

\begin{figure}[H]
\centering
\includegraphics[width=0.9\linewidth]{pic/1_polar_dark/2_polar_channels.png}
\caption{暗弱光照条件下的斯托克斯矢量与偏振通道对比}
\label{fig:polar_channels}
\end{figure}

图\ref{fig:polar_channels}为3行6列的子图矩阵。每一行对应一种光照条件：第1行为正常光照（94--205 lux），第2行为较暗光照（2.6--18.6 lux），第3行为极暗光照（0.9--4.8 lux）。从左到右6列依次为：可见光强度$S_0$、斯托克斯分量$S_1$（取绝对值后放大30倍显示）、斯托克斯分量$S_2$（取绝对值后放大30倍显示）、偏振度$P_0(DoP)$、偏振角正弦$P_1(\sin AoP)$、偏振角余弦$P_2(\cos AoP)$。

从图\ref{fig:polar_channels}可见，正常光照条件下，斯托克斯分量$S_1, S_2$富含场景结构信息，由此导出的偏振通道$P_0, P_1, P_2$同样保留了丰富的结构纹理，这些信息均可作为视觉导航的有效特征。随着光照减弱，可见光图像$S_0$的纹理特征迅速退化，而偏振通道在极暗条件下（0.9--4.8 lux）仍可辨识出场景结构。

上述现象的物理根源在于偏振信息与强度信息承载了不同类型的场景对比度。以下从偏振信号的物理来源、偏振参量的光照不变性、强度图像的退化机制、以及传感器共模噪声抑制四个层面予以分析。

在室内环境中，偏振信号的主要来源是物体表面的菲涅耳（Fresnel）反射。对于介质表面的镜面反射，偏振度$DoP$由菲涅耳反射率$R_s$与$R_p$的比值决定：$DoP = |R_s - R_p| / (R_s + R_p)$。其中$R_s$和$R_p$仅取决于光线的入射角、材料折射率及表面粗糙度等几何与材质属性\cite{born1999principles}，与入射光的绝对强度无关。偏振角$AoP$则由入射平面方向决定，编码了表面法线的方向信息。因此，$DoP$与$AoP$所承载的是场景的"几何/材质对比度"（geometry/material contrast），而非普通强度图像$S_0$所依赖的"光照对比度"（illumination contrast）。

上述物理来源赋予偏振通道对光照强度的固有不变性。即：当环境光照整体减弱时，四个偏振方向的原始光强$I_0, I_{45}, I_{90}, I_{135}$等比衰减，Stokes参量$S_0, S_1, S_2$均按相同因子缩放。偏振度$DoP = \sqrt{S_1^2 + S_2^2} / S_0$为比值量，分子与分母中的缩放因子相互抵消，因此理论上对光照强度变化具有一阶不变性。类似地，偏振角$AoP = \frac{1}{2}\arctan(S_2/S_1)$仅依赖于$S_2$与$S_1$的比值，同样不受光照等比衰减的影响。不过，该不变性在极低信噪比条件下会逐渐退化：$S_0 \to 0$时$DoP$出现除零不稳定，$S_1, S_2 \approx 0$时$AoP$趋于随机化。但在信噪比尚可维持的范围内，偏振参量对光照强度的不敏感性已足以为视觉特征提供稳定的结构信息。

上述偏振信息与$S_0$信息来源的根本差异，决定了两者在暗光下不同的退化模式。$S_0$的纹理对比度来源于光照非均匀性——阴影、照明方向性造成的亮度梯度——与表面反射率的共同作用。当环境光照降至数lux量级，光子散粒噪声方差$\sigma_{shot}^2 \propto I$相对于信号幅度不再可忽略，依赖光照梯度的纹理信息首先被噪声淹没。而偏振通道$P_0, P_1, P_2$承载的是与光照无关的几何/材质对比度，其结构信息在$S_0$失效后仍可持续存在。简言之，$S_0$编码光照主导的对比度，偏振通道编码几何/材质主导的对比度；暗光条件下前者率先失效，后者得以保留。

以上从信号来源与物理成分的角度解释了偏振通道结构信息在暗光下得以保持的原因，但尚需澄清一个信号处理层面的问题：$S_1 = I_0 - I_{90}$为差分运算，经典误差传播理论给出$\sigma_{S_1}^2 = \sigma_{I_0}^2 + \sigma_{I_{90}}^2$，形式上噪声方差应被放大，与实验观测存在表观矛盾。

该矛盾的关键在于，上述误差传播公式隐含了$I_0$与$I_{90}$噪声相互独立的假设。在微阵列偏振相机（DoFP）中，同一超像素的四个子像素物理上紧密相邻、共享同一微透镜和读出电路，其暗电流噪声、读出噪声与ADC量化噪声中存在显著的共模分量\cite{perkins2010signal,bai2022noise}。将子像素光强建模为$I_k = s_k + n_c + n_k$，其中$n_c$为四个子像素共有的共模噪声，$n_k$为各子像素的独立噪声，则$S_1 = (s_0 - s_{90}) + (n_0 - n_{90})$——共模分量$n_c$在差分运算中被完全抵消，仅独立噪声分量残留。从信号处理的角度，该差分过程可类比为锁相放大器（lock-in amplifier）的相敏检波：$I_0$和$I_{90}$分别测量了偏振调制信号在两个正交相位上的响应，差分运算在抑制共模直流背景的同时提取偏振调制幅度，因此有效信噪比并未如独立噪声假设所预期的那样严重恶化。

综上，偏振信息在暗光下的稳健性根植于一个核心机制：偏振通道通过差分提取与比值归一化，系统性地抑制了与光照强度相关的信号成分，保留与几何结构和材料属性相关的信号成分。这使得偏振特征在光照强度跨越约两个数量级（$\sim$200 lux至$\sim$1 lux）的范围内，仍能提供可用的结构信息。

然而，上述物理机制并不意味着偏振通道在暗光下不受噪声影响。从图\ref{fig:polar_channels}中亦可观察到，随光照减弱，偏振通道中的图像噪声显著增大；在极暗条件下，噪声已严重劣化图像质量。其原因是，在暗光条件下$I_0, I_{45}, I_{90}, I_{135}$等原始信号微弱，信噪比较低，而传感器暗电流等固有噪声源无法物理消除。这些噪声经偏振通道提取（式~\ref{eq:stokes}至\ref{eq:polar-channel-2}）过程中的差分与非线性运算传播后进一步放大，严重干扰后续的特征提取与匹配。因此，有必要通过图像预处理算法对其进行抑制。


\textbf{(2) 偏振通道预处理方法}


暗光条件下，图像传感器内部暗电流等因素会造成噪声。根据图\ref{fig:polar_channels}可知，这些噪声近似于图像椒盐噪声。因此，考虑通过图像滤波算法对其进行去噪预处理。

然而，由于我们最终目标是解决极端暗光环境下的视觉导航问题，而视觉导航依赖于清晰的视觉特征追踪。如果使用常见的高斯滤波等算法进行去噪，图像边缘将变得模糊，从而使后续特征提取算法很难提取清晰可用的特征。

为抑制上述噪声，同时尽可能保留图像边缘以供后续特征提取，需要分析不同滤波算法的适用性。
\begin{itemize}
    \item 中值滤波通过邻域像素排序取中值，对椒盐噪声和脉冲型噪声具有天然鲁棒性，计算开销小，但缺乏对图像边缘的显式保护，多次迭代可能导致角点模糊。
    \item 双边滤波\cite{tomasi1998bilateral}结合空间域和颜色域高斯加权，可在平滑噪声的同时保留强边缘；然而，偏振通道中的残余噪声在灰度分布上与真实边缘可能具有相近的方差，导致双边滤波在暗光下将噪声误判为边缘而保留。
    \item 非局部均值去噪\cite{buades2005nlm}(Non-Local Means, NLM)利用图像块自相似性，理论上可获得优异的噪声抑制，但其计算复杂度极高，需要对大范围搜索窗口进行相似性匹配，难以满足视觉导航的实时性约束。
    \item 导向滤波\cite{he2013guided}以另一幅引导图像的局部线性模型为基础进行滤波，当以$S_0$强度图为引导时，可利用$S_0$中仍保留的边缘结构来指导偏振通道的平滑过程，在去噪的同时显式保持边缘。
\end{itemize}

基于上述分析，本文选取这四种具有代表性的算法进行对比实验。为探究最佳的偏振通道图像预处理算法，我们选用上述4种图像滤波算法，对3个偏振通道的图像进行滤波测试。下列测试中，各算法使用的参数如表\ref{tab:filter_params}所示。这些参数都是多次实验中选取的效果最佳的参数。

\begin{table}[htbp]
\centering
\caption{偏振通道预处理滤波算法参数设置}
\label{tab:filter_params}
\begin{tabular}{@{}lccc@{}}
\toprule
滤波算法 & 参数 & 取值 & 说明 \\
\midrule
中值滤波 (Median)     & 核大小        & $3\times 3$  & 执行3轮滤波 \\
\addlinespace
导向滤波 (Guided)     & 窗口半径 $r$  & 8            & 以 $S_0$ 强度图像为引导 \\*
                     & 正则化 $\varepsilon$ & 0.0001 & \\
\addlinespace
双边滤波 (Bilateral)  & 邻域直径 $d$  & 9           & \\*
                     & $\sigma_{\text{color}}$ & 200  & 颜色空间标准差 \\*
                     & $\sigma_{\text{space}}$ & 30   & 空间域标准差 \\
\addlinespace
非局部均值去噪 (NLM)  & 滤波强度 $h$  & 50           & \\*
                     & 模板窗口       & $5\times 5$  & \\*
                     & 搜索窗口       & $21\times 21$ & \\
\bottomrule
\end{tabular}
\end{table}


实验条件与本节第(1)小节相同，采集 3 组偏振图像，每组 200 帧，帧率 7FPS。

采用上述3组数据，分别对上述4种算法进行测试。选取部分数据的滤波效果进行展示，如图\ref{fig:polar_filter_1}至图\ref{fig:polar_filter_3}所示。

\begin{figure}[htbp]
\centering
\includegraphics[width=0.9\linewidth]{pic/1_polar_dark/3_filter_1.png}
\caption{正常光照条件(94-205 lux)下偏振通道滤波算法对比}
\label{fig:polar_filter_1}
\end{figure}

\begin{figure}[htbp]
\centering
\includegraphics[width=0.9\linewidth]{pic/1_polar_dark/3_filter_2.png}
\caption{较暗光照条件(2.6-18.6 lux)下偏振通道滤波算法对比}
\label{fig:polar_filter_2}
\end{figure}

\begin{figure}[htbp]
\centering
\includegraphics[width=0.9\linewidth]{pic/1_polar_dark/3_filter_3.png}
\caption{极暗光照条件(0.9-4.8 lux)下偏振通道滤波算法对比}
\label{fig:polar_filter_3}
\end{figure}

图\ref{fig:polar_filter_1}至图\ref{fig:polar_filter_3}均为3行5列的子图矩阵。每一行对应一个偏振通道：第1行为$P_0(DoP)$，第2行为$P_1(\sin AoP)$，第3行为$P_2(\cos AoP)$。从左到右5列依次为：无滤波(Raw)、中值滤波(Median)、导向滤波(Guided)、双边滤波(Bilateral)、非局部均值去噪(NLM)。三个图分别对应不同光照条件：图\ref{fig:polar_filter_1}为正常光照条件，图\ref{fig:polar_filter_2}为较暗光照条件，图\ref{fig:polar_filter_3}为极暗光照条件。

从上述对比可以看到，在对比的4种滤波算法中，$S_0$引导的导向滤波去噪效果最好，在去噪的同时，仍保留了部分锐利的边缘信息。其余算法对暗光情况下偏振通道噪声的抑制均不如$S_0$引导的导向滤波。


由于后续需要基于各偏振通道提取视觉特征点，并在不同帧特征点间进行匹配，使用特征点提取与匹配方法计算量化指标，评估不同滤波算法。

具体来说，选用GFTT(Good Features to Track, Shi \& Tomasi, CVPR 1994)\cite{shi1994good}特征点作为各通道特征点检测方法，选用L-K光流法\cite{lucas1981lucas}对相邻帧检测的特征点进行匹配，匹配后使用本质矩阵RANSAC\cite{fischler1981ransac}方法检测内点。统计每帧检测出的特征点数量，与每相邻两帧特征点匹配的内点数量，以帧平均内点数量和帧平均内点比例作为图像预处理方法的量化评价指标。同时，为满足视觉导航实时性要求，统计帧平均处理时间，即每帧图像偏振通道提取-偏振通道预处理-特征点检测-特征点匹配-RANSAC的总时间，也作为图像预处理方法的量化评价指标。


\begin{figure}[H]
\centering
\includegraphics[width=0.9\linewidth]{pic/1_polar_dark/4_filter_metrics.png}
\caption{不同光照条件下偏振通道滤波算法特征提取指标对比}
\label{fig:polar_filter_4}
\end{figure}

使用上述评价指标，在搭载英特尔i7-12700H处理器的计算机上运行测试，在3个偏振通道上测试上述4种滤波方法的效果。各算法量化指标如图\ref{fig:polar_filter_4}和表\ref{tab:filter_metrics}所示。

\begin{table*}[t]
\centering
\caption{不同滤波方法与偏振通道组合的特征跟踪量化结果}
\label{tab:filter_metrics}
\footnotesize
\setlength{\tabcolsep}{3.5pt}
\begin{tabular}{llccccccccc}
\toprule
& & \multicolumn{3}{c}{\bfseries 明亮光照 (94--205~lux)} & \multicolumn{3}{c}{\bfseries 中等光照 (2.6--18.6~lux)} & \multicolumn{3}{c}{\bfseries 暗光 (0.9--4.8~lux)} \\
\cmidrule(lr){3-5} \cmidrule(lr){6-8} \cmidrule(lr){9-11}
\multirow{-2}{*}{\bfseries 滤波方法} & \multirow{-2}{*}{\bfseries 通道} & 内点数$\uparrow$ & 内点比例$\uparrow$ & 时间/ms$\downarrow$ & 内点数$\uparrow$ & 内点比例$\uparrow$ & 时间/ms$\downarrow$ & 内点数$\uparrow$ & 内点比例$\uparrow$ & 时间/ms$\downarrow$ \\
\midrule

% ---- Raw ----
\multirow{3}{*}{Raw}
    & $P_0$ (DoP)     & \textbf{250.9} & 0.868 & 4.7  & \textbf{98.2} & 0.608 & \textbf{7.3}  & \textbf{43.8} & 0.360 & 11.9 \\
    & $P_1$ (sinAoP)  & \textbf{277.7} & 0.822 & 5.3  & \textbf{46.6} & 0.338 & 12.0 & \textbf{32.2} & 0.244 & 12.3 \\
    & $P_2$ (cosAoP)  & \textbf{262.8} & 0.802 & 5.1  & \textbf{61.2} & 0.422 & 11.5 & \textbf{41.1} & 0.303 & 12.2 \\
\cmidrule{2-11}

% ---- Median ----
\multirow{3}{*}{中值滤波}
    & $P_0$ (DoP)     & 172.9 & 0.876 & \textbf{4.1}  & 83.7  & 0.574 & 7.4  & 32.7  & 0.359 & 11.1 \\
    & $P_1$ (sinAoP)  & 265.5 & 0.834 & \textbf{5.0}  & 39.6  & 0.372 & 11.7 & 23.9  & 0.262 & 11.9 \\
    & $P_2$ (cosAoP)  & 246.3 & 0.802 & \textbf{4.9}  & 57.5  & 0.463 & 10.7 & 30.8  & 0.316 & 11.8 \\
\cmidrule{2-11}

% ---- Guided ----
\multirow{3}{*}{导向滤波}
    & $P_0$ (DoP)     & 152.7 & \textbf{0.910} & 8.7  & 62.8  & \textbf{0.710} & 9.0  & 37.2  & \textbf{0.622} & \textbf{10.0} \\
    & $P_1$ (sinAoP)  & 188.4 & \textbf{0.839} & 9.2  & 24.4  & \textbf{0.653} & \textbf{9.8}  & 11.5  & \textbf{0.549} & \textbf{11.4} \\
    & $P_2$ (cosAoP)  & 213.2 & \textbf{0.883} & 9.1  & 53.2  & \textbf{0.781} & \textbf{8.7}  & 20.6  & \textbf{0.650} & \textbf{9.7} \\
\cmidrule{2-11}

% ---- Bilateral ----
\multirow{3}{*}{双边滤波}
    & $P_0$ (DoP)     & 171.0 & 0.842 & 7.4  & 78.9  & 0.588 & 10.3 & 26.8  & 0.406 & 13.9 \\
    & $P_1$ (sinAoP)  & 244.4 & 0.810 & 8.3  & 28.4  & 0.391 & 14.8 & 19.3  & 0.312 & 15.1 \\
    & $P_2$ (cosAoP)  & 232.1 & 0.784 & 8.2  & 48.3  & 0.489 & 13.4 & 25.1  & 0.359 & 15.0 \\
\cmidrule{2-11}

% ---- NLM ----
\multirow{3}{*}{NLM}
    & $P_0$ (DoP)     & 58.5  & 0.876 & 136.1 & 56.4  & 0.583 & 146.8 & 33.3  & 0.362 & 147.5 \\
    & $P_1$ (sinAoP)  & 235.7 & 0.839 & 141.7 & 41.5  & 0.348 & 155.0 & 29.8  & 0.248 & 151.3 \\
    & $P_2$ (cosAoP)  & 207.0 & 0.795 & 139.7 & 60.2  & 0.443 & 150.7 & 38.6  & 0.304 & 147.0 \\

\midrule
\end{tabular}

\vspace{4pt}
\begin{minipage}{\textwidth}
\footnotesize
\textbf{注：}
内点数（Avg Inlier）为帧平均 RANSAC 基础矩阵内点数，越高越好，反映后端可用的几何一致特征匹配数量.
内点比例（Inlier Ratio）为内点数与匹配数之比，反映滤波方法的匹配精度.
时间为一帧内偏振通道提取-偏振通道预处理-特征点检测-特征点匹配-RANSAC的全流程耗时.
\end{minipage}
\end{table*}

图\ref{fig:polar_filter_4}为3行3列的子图矩阵。每一行对应一种光照条件：第1行为明亮光照（94--205 lux），第2行为中等光照（2.6--18.6 lux），第3行为暗光（0.9--4.8 lux）。从左到右3列依次为：内点数（Avg Inlier Count）、内点比例（Inlier Ratio）、单帧处理时间（Avg Time，对数坐标）。每个子图中，横坐标为5种滤波方法（Raw、Median、Guided、Bilateral、NLM），每组包含3个彩色柱状条，分别对应3个偏振通道：蓝色为$DoP(P_0)$，橙色为$\sin AoP(P_1)$，绿色为$\cos AoP(P_2)$。

从图\ref{fig:polar_filter_4}和表\ref{tab:filter_metrics}可以看到，随环境光照强度变弱，各算法检测出的特征点内点数量均下降。需注意的是，导向滤波在明亮光照下的内点数（如$P_0$通道152.7）低于无滤波方法（250.9），这并非说明导向滤波效果更差，而是由于导向滤波在去除噪声的同时过滤掉了大量不可靠的候选特征点，使得剩余特征点的几何一致性显著提高，最终体现为最高的内点比例（0.910）。在视觉SLAM中，非内点的特征点会干扰优化算法，可能导致跟踪丢失和优化失败等，因此\textbf{内点比例}是比内点数量更核心的指标。

此外，NLM虽然也能获得较高的内点比例，但其单帧处理时间超过130ms，远超视觉导航系统的实时性约束（通常要求每帧总耗时低于33ms），故在实际系统中不具备可用性。综合来看，导向滤波在各类滤波方法中优势明显，且单通道单帧处理时间均在10ms左右，满足视觉导航系统的实时性要求。

综合上述可视化分析与定量分析结果，选取导向滤波作为偏振通道预处理方法。


\textbf{(3) 偏振特征提取方法}

通过偏振通道提取与偏振通道预处理，就获得了可用于暗弱环境下视觉导航的偏振通道图像。要想从这些图像中获取导航信息，需要提取特征点并进行匹配。

图像特征点提取领域有许多经过长期实践检验和验证的方法，被广泛应用于多视角几何和视觉导航领域，如FAST和GFTT等。除了传统方法，近年来SuperPoint等基于深度神经网络的方法也在许多场景下取得了最佳效果。

然而，偏振通道不同于普通RGB相机获取的可见光图像。偏振通道的数据分布，尤其是经过上述滤波预处理后的残余噪声分布，不同于可见光图像。因此，仍采用定量指标测试的方法，对比不同特征点提取算法在预处理后的偏振通道上提取和匹配特征点的表现，从而选取最合适的特征提取方法。

具体来说，选用FAST\cite{rosten2006fast}、GFTT\cite{shi1994good}、SuperPoint\cite{deTone2018superpoint}三种特征点作为待评估的特征点检测方法。


\begin{figure}[htbp]
\centering
\includegraphics[width=0.8\linewidth]{pic/1_polar_dark/5_fp_vis_1.png}
\caption{正常光照条件(94-205 lux)下偏振特征提取算法对比}
\label{fig:polar_fp_1}
\end{figure}

仍使用上一小节滤波方法对比测试中所用数据作为评估测试数据，预处理方法统一使用上节选定的导向滤波。各组数据采集条件与上节相同，每组数据共包含200张单目偏振相机连续采集的偏振图像，采集帧率为7FPS。

\begin{figure}[htbp]
\centering
\includegraphics[width=0.8\linewidth]{pic/1_polar_dark/5_fp_vis_2.png}
\caption{较暗光照条件(2.6-18.6 lux)下偏振特征提取算法对比}
\label{fig:polar_fp_2}
\end{figure}


采用上述3组数据，分别对上述4种算法进行测试。选取部分数据的特征提取效果进行展示，如图\ref{fig:polar_fp_1}至图\ref{fig:polar_fp_3}所示。


\begin{figure}[htbp]
\centering
\includegraphics[width=0.8\linewidth]{pic/1_polar_dark/5_fp_vis_3.png}
\caption{极暗光照条件(0.9-4.8 lux)下偏振特征提取算法对比}
\label{fig:polar_fp_3}
\end{figure}

图\ref{fig:polar_fp_1}至图\ref{fig:polar_fp_3}均为3行3列的子图矩阵。每一行对应一个偏振通道：第1行为$P_0(DoP)$，第2行为$P_1(\sin AoP)$，第3行为$P_2(\cos AoP)$。从左到右3列依次为：FAST特征检测器、GFTT特征检测器、SuperPoint特征检测器。绿色散点表示各算法在当前通道上检测到的特征点位置。图\ref{fig:polar_fp_1}为正常光照条件，图\ref{fig:polar_fp_2}为较暗光照条件，图\ref{fig:polar_fp_3}为极暗光照条件。

从上述对比可以看到，在对比的3种特征检测算法中，FAST算法提取的特征点较为集中，SuperPoint提取的特征点较为分散，而GFTT提取的特征点较为均衡。在暗光场景下，GFTT可以保留相对较多的图像特征。

使用特征点提取与匹配方法计算量化指标，量化评估不同特征点提取算法。具体来说，仍选用L-K光流法\cite{lucas1981lucas}对相邻帧检测的特征点进行匹配，匹配后仍使用本质矩阵RANSAC\cite{fischler1981ransac}方法检测内点。统计每帧检测出的特征点数量，与每相邻两帧特征点匹配的内点数量，以帧平均内点数量和帧平均内点比例作为特征点检测方法的量化评价指标。同时，为满足视觉导航实时性要求，统计帧平均处理时间，即每帧图像偏振通道提取-偏振通道预处理-特征点检测-特征点匹配-RANSAC的总时间，也作为图像预处理方法的量化评价指标。

\begin{figure}[H]
\centering
\includegraphics[width=0.9\linewidth]{pic/1_polar_dark/6_fp_metric.png}
\caption{不同光照条件下偏振特征提取算法特征提取指标对比}
\label{fig:polar_fp_4}
\end{figure}

使用上述评价指标，在搭载英特尔i7-12700H处理器的计算机上运行测试，在3个偏振通道上测试上述4种滤波方法的效果。各算法量化指标如图\ref{fig:polar_fp_4}和表\ref{tab:fp_metrics}所示。


\begin{table*}[t]
\centering
\caption{不同特征检测方法在偏振通道上的特征跟踪量化结果}
\label{tab:fp_metrics}
\footnotesize
\setlength{\tabcolsep}{3.5pt}
\begin{tabular}{llccccccccc}
\toprule
& & \multicolumn{3}{c}{\bfseries 明亮光照 (94--205~lux)} & \multicolumn{3}{c}{\bfseries 中等光照 (2.6--18.6~lux)} & \multicolumn{3}{c}{\bfseries 暗光 (0.9--4.8~lux)} \\
\cmidrule(lr){3-5} \cmidrule(lr){6-8} \cmidrule(lr){9-11}
\multirow{-2}{*}{\bfseries 检测方法} & \multirow{-2}{*}{\bfseries 通道} & 内点数$\uparrow$ & 内点比例$\uparrow$ & 时间/ms$\downarrow$ & 内点数$\uparrow$ & 内点比例$\uparrow$ & 时间/ms$\downarrow$ & 内点数$\uparrow$ & 内点比例$\uparrow$ & 时间/ms$\downarrow$ \\
\midrule

% ---- FAST ----
\multirow{3}{*}{FAST}
    & $P_0$ (DoP)     & 73.4  & \textbf{0.943} & \textbf{1.0}  & 1.5   & 0.370 & \textbf{0.8}  & 0.0   & 0.036 & \textbf{0.4} \\
    & $P_1$ (sinAoP)  & 121.9 & \textbf{0.910} & \textbf{1.1}  & 1.0   & 0.299 & \textbf{0.7}  & 0.0   & 0     & \textbf{0.3} \\
    & $P_2$ (cosAoP)  & \textbf{304.6} & \textbf{0.911} & \textbf{1.6}  & 20.1  & \textbf{0.878} & \textbf{1.1}  & 1.9   & 0.465 & \textbf{0.9} \\
\cmidrule{2-11}

% ---- GFTT ----
\multirow{3}{*}{GFTT}
    & $P_0$ (DoP)     & \textbf{152.7} & 0.909 & 3.8  & \textbf{62.8} & 0.710 & 4.2  & \textbf{37.2} & \textbf{0.622} & 5.3 \\
    & $P_1$ (sinAoP)  & \textbf{188.4} & 0.839 & 4.0  & \textbf{24.4} & \textbf{0.653} & 4.7  & \textbf{11.5} & \textbf{0.549} & 6.3 \\
    & $P_2$ (cosAoP)  & 213.2 & 0.882 & 3.9  & \textbf{53.2} & 0.781 & 3.5  & \textbf{20.6} & \textbf{0.650} & 4.6 \\
\cmidrule{2-11}

% ---- SuperPoint ----
\multirow{3}{*}{SuperPoint}
    & $P_0$ (DoP)     & 25.5  & 0.898 & 9.9  & 12.9  & \textbf{0.762} & 10.6 & 5.0   & 0.555 & 10.2 \\
    & $P_1$ (sinAoP)  & 23.0  & 0.847 & 9.9  & 2.4   & 0.361 & 10.6 & 0.1   & 0.032 & 10.2 \\
    & $P_2$ (cosAoP)  & 25.9  & 0.885 & 9.9  & 10.3  & 0.720 & 10.6 & 2.7   & 0.413 & 10.2 \\

\midrule
\end{tabular}

\vspace{4pt}
\begin{minipage}{\textwidth}
\footnotesize
\textbf{注：}
内点数（Avg Inlier）为帧平均 RANSAC 基础矩阵内点数，越高越好，反映后端可用的几何一致特征匹配数量.
内点比例（Inlier Ratio）为内点数与匹配数之比，反映检测方法的匹配精度.
时间为一帧内偏振通道提取-导向滤波预处理-特征点检测-特征点匹配-RANSAC的全流程耗时.
\end{minipage}
\end{table*}



图\ref{fig:polar_fp_4}为3行3列的子图矩阵。每一行对应一种光照条件：第1行为明亮光照（94--205 lux），第2行为中等光照（2.6--18.6 lux），第3行为暗光（0.9--4.8 lux）。从左到右3列依次为：内点数（Avg Inlier Count）、内点比例（Inlier Ratio）、单帧处理时间（Avg Time，对数坐标）。每个子图中，横坐标为3种特征检测方法（FAST、GFTT、SuperPoint），每组包含3个彩色柱状条，分别对应3个偏振通道：蓝色为$DoP(P_0)$，橙色为$\sin AoP(P_1)$，绿色为$\cos AoP(P_2)$。

从图\ref{fig:polar_fp_4}和表\ref{tab:fp_metrics}可以看到，随环境光照强度变弱，各算法检测出的特征点内点数量均下降，但GFTT方法检测特征点数量相对最多。

值得注意的是，SuperPoint作为基于深度学习的特征检测器，在偏振通道上的性能显著低于其在普通RGB图像上的典型表现。这主要是因为我们使用的是SuperPoint官方提供的预训练模型(MagicLeap官方开源实现\cite{deTone2018superpoint}，2018年发布，使用COCO和HPatches数据集训练的PyTorch权重)，其训练数据均为自然场景的RGB灰度图像；而偏振通道（尤其是$DoP$和$AoP$分量）的物理含义、噪声分布和对比度特性与RGB图像存在本质差异，导致预训练模型无法有效提取偏振图像中的稳定特征。该结果也表明，暗光偏振视觉导航不能简单套用常规RGB图像上的深度学习方法，而需要针对偏振数据特性设计或重新训练专用检测器。

在内点比例方面，对于滤波后的各偏振通道，即使极暗光照条件下，GFTT也可以维持60\%左右的内点比例，说明其检测出的特征点质量较高。在处理时间方面，GFTT在各光照条件下的各偏振通道处理时间均为4ms左右，满足视觉导航系统的实时性要求。

综合上述可视化分析与定量分析结果，选取GFTT方法作为偏振通道特征点检测方法。
